% Part: second-order-logic
% Chapter: sol-and-sets
% Section: introduction

\documentclass[../../../include/open-logic-section]{subfiles}

\begin{document}

\olfileid{sol}{set}{int}

\olsection{परिचय}

द्वितीय-क्रम तर्कशास्त्रात प्रांताच्या उपसंचांवर तसेच
फलनांवर संख्यापन करता येते; म्हणून त्यात संचसिद्धान्ताचा
काही भाग तरी मांडता येईल अशी अपेक्षा आहे. येथे
‘मांडणे’ याचा अर्थ संचसिद्धान्तीय गुणधर्म आणि विधाने
द्वितीय-क्रम तर्कशास्त्रात व्यक्त करता येणे, तेही
संचांसाठी कोणतीही खास अतार्किक चिन्हे (उदा.,
संचसिद्धान्तातील सदस्यत्वाचे !!{predicate}) न वापरता.
उदाहरणार्थ, संचांचे संयोग आणि छेद तसेच उपसंच-संबंध
परिभाषित करता येतात; शिवाय संचांच्या आकारांची तुलना
करता येते आणि कँटरच्या प्रमेयासारखे निष्कर्षही मांडता येतात.

\end{document}
